% AUTO-GENERATED by Github-Branch/tools/build_unified_source.ps1.
% Do not edit this file directly: update the transformation manifest instead.
\chapter{Measure, Probability and Stochastic Processes}
\chapterstatus{Edited reuse from Team 1. Mathematical and source-to-claim review remains open; figures are withheld pending the figure ledger.}
\sourceassetnote
\section{Measure theory and probability introduction}
\subsection{Measurable Spaces}

A measurable space is a pair consisting of a set and a $\sigma$-algebra.
A particularly important example is a Borel space, where the $\sigma$-algebra is the Borel $\sigma$-algebra of a topological space.

\par
A measurable space captures and generalises intuitive notions such as length, area, and volume with a set $X$ of ``points'' in the space. However, regions of the space are the elements of the $\sigma$-algebra, since intuitive measures are not usually defined for individual points. The algebra also captures relationships between regions: a region can be defined as an intersection of other regions, a union of other regions, or the space with the exception of another region.

\par
\textbf{Definition.} Consider a set $X$ and a $\sigma$-algebra $\mathcal{F}$ on $X$. Then the tuple $(X, \mathcal{F})$ is called a \textit{measurable space}. The elements of $\mathcal{F}$ are called \textit{measurable sets} within the measurable space.

\par
Note that, in contrast to a \textit{measure space}, no measure is needed for a measurable space.

\par
\textbf{Example.} Consider the set $X = \{1, 2, 3\}$. One possible $\sigma$-algebra would be 
\[
\mathcal{F}_1 = \{ X, \emptyset \}.
\]
Then $(X, \mathcal{F}_1)$ is a measurable space. Another possible $\sigma$-algebra is the power set of $X$:
\[
\mathcal{F}_2 = \mathcal{P}(X).
\]
With this, a second measurable space on the set $X$ is given by $(X, \mathcal{F}_2)$.

\par
\textbf{Common measurable spaces.}  
If $X$ is finite or countably infinite, the $\sigma$-algebra is most often the power set on $X$, so $\mathcal{F} = \mathcal{P}(X)$. This leads to the measurable space $(X, \mathcal{P}(X))$.

\par
If $X$ is a topological space, the $\sigma$-algebra is most commonly the \textit{Borel $\sigma$-algebra} $\mathcal{B}$, so $\mathcal{F} = \mathcal{B}(X)$. This leads to the measurable space $(X, \mathcal{B}(X))$, which is common for all topological spaces such as the real numbers $\mathbb{R}$.

\par
\textbf{Probability spaces.}  
A probability space is a measurable space equipped with a measure that is normalized. Formally, a probability space is a triple $(\Omega, \mathcal{F}, \mathbb{P})$, where $(\Omega, \mathcal{F})$ is a measurable space and $\mathbb{P} : \mathcal{F} \to [0,1]$ is a measure such that $\mathbb{P}(\Omega) = 1$. The set $\Omega$ is called the sample space, and elements of $\mathcal{F}$ are called events.
\subsection{Sigma-Algebras}

In mathematical analysis and probability theory, a \textbf{$\sigma$-algebra} (``sigma algebra'') is part of the formalism for defining sets that can be measured. In calculus and analysis, for example, $\sigma$-algebras are used to define the concept of sets with area or volume. In probability theory, they are used to define events with a well-defined probability. In this way, $\sigma$-algebras help to formalize the notion of size.

\par
In formal terms, a $\sigma$-algebra (also called a \textbf{$\sigma$-field}, where the $\sigma$ comes from the German word \textit{``Summe''}, meaning ``sum'') on a set $X$ is a nonempty collection $\Sigma \subseteq \mathcal{P}(X)$ of subsets of $X$ such that:
\begin{itemize}
    \item $X \in \Sigma$,
    \item if $A \in \Sigma$, then its complement $A^c := X \setminus A \in \Sigma$,
    \item if $(A_n)_{n \in \mathbb{N}} \subseteq \Sigma$, then $\bigcup_{n=1}^{\infty} A_n \in \Sigma$.
\end{itemize}
From these properties it follows that $\Sigma$ is also closed under countable intersections.

\par
The ordered pair $(X, \Sigma)$ is called a \textit{measurable space}.

\par
The set $X$ is understood to be an ambient space (such as the two-dimensional plane or the set of outcomes when rolling a six-sided dice, $\{1,2,3,4,5,6\}$), and the collection $\Sigma$ is a choice of subsets declared to have a well-defined size. The closure requirements for $\sigma$-algebras are designed to capture our intuitive ideas about how sizes combine.

\par
The definition of a $\sigma$-algebra resembles other mathematical structures such as a \textit{topology} or a \textit{set algebra}.
\paragraph{Examples of $\sigma$-Algebras}

If $X = \{ a, b, c, d \}$, one possible $\sigma$-algebra on $X$ is
\[
\Sigma = \{ \varnothing, \{a,b\}, \{c,d\}, \{a,b,c,d\} \}.
\]

\par
If $\{A_1, A_2, A_3, \ldots\}$ is a countable partition of $X$, then the collection of all unions of sets in the partition (including the empty set) is a $\sigma$-algebra.

\par
A more useful example is the set of subsets of the real line $\mathbb{R}$ formed by starting with all open intervals and closing under countable unions, countable intersections, and complements. This construction is known as the \textit{Borel $\sigma$-algebra}.
\subsection{Measure}

A \textit{measure} on $X$ is a function
\[
\mu : \mathcal{F} \to [0, +\infty]
\]
defined on a $\sigma$-algebra $\mathcal{F}$ such that:
\begin{itemize}
    \item $\mu(\emptyset) = 0$,
    \item (countable additivity) if $(A_n)_{n \in \mathbb{N}}$ are pairwise disjoint sets in $\mathcal{F}$, then
    \[
    \mu\left( \bigcup_{n=1}^{\infty} A_n \right) = \sum_{n=1}^{\infty} \mu(A_n).
    \]
\end{itemize}

\par
This formalizes the idea of assigning a consistent notion of ``size'' to sets.

\par
Ideally, one would like to assign a size to every subset of $X$, but in many natural settings this is not possible. For example, the \textit{axiom of choice} implies that there exist subsets of $\mathbb{R}$ (such as Vitali sets) that are not measurable.

\par
For this reason, one considers instead a smaller collection of privileged subsets of $X$, called the \textit{measurable sets}, which form a $\sigma$-algebra.

\section{Random variables, Lebesgue integrals and convergence theorems}
\subsection{Random variables}

A \textbf{random variable} $X$ is a measurable function
\[
X : \Omega \to E
\]
from a sample space $\Omega$ (the set of possible outcomes) to a measurable space $(E, \mathcal{E})$, where typically $E = \mathbb{R}$ or $\mathbb{R}^n$ equipped with the Borel $\sigma$-algebra.

\par
The technical, axiomatic definition requires the sample space $\Omega$ to belong to a \textit{probability triple} $(\Omega, \mathcal{F}, P)$, where $\mathcal{F}$ is a $\sigma$-algebra of events and $P$ is a probability measure.

\par
The measurability of $X$ means that for every measurable set $S \in \mathcal{E}$,
\[
X^{-1}(S) := \{ \omega \in \Omega \mid X(\omega) \in S \} \in \mathcal{F}.
\]

\par
The probability that $X$ takes on a value in a measurable set $S \subseteq E$ is therefore written as
\[
P(X \in S) = P(X^{-1}(S)).
\]
\subsection{Density function, expectation, variance}

A real-valued random variable $X$ is described by means of its cumulative distribution function (CDF),
\[
F_X(x) := \mathbb{P}[X \leq x].
\]

\par
If $X$ has an \textit{absolutely continuous} distribution, then it admits a probability density function (PDF) $f_X$ such that
\[
F_X(x) = \int_{-\infty}^{x} f_X(t)\,dt, \qquad \text{and} \qquad f_X(x) = \frac{dF_X(x)}{dx}.
\]
In general, not every random variable admits a density: discrete or mixed distributions do not have a PDF in this sense.

\par
The expected value of a random variable $X$ is defined in general (measure-theoretically) as the Lebesgue integral
\[
\mathbb{E}[X] = \int_{\Omega} X(\omega)\, dP(\omega),
\]
provided the integral $\int_{\Omega} |X|\, dP < \infty$.

\par
In the absolutely continuous case with density $f_X$, this reduces to
\[
\mathbb{E}[X] = \int_{-\infty}^{+\infty} x f_X(x)\, dx.
\]

\par
The variance of $X$, $\mathrm{Var}[X]$, is defined as
\[
\mathrm{Var}[X] = \int_{\Omega} (X - \mathbb{E}[X])^2\, dP,
\]
whenever the integral exists. In the case of a density,
\[
\mathrm{Var}[X] = \int_{-\infty}^{+\infty} (x - \mathbb{E}[X])^2 f_X(x)\, dx.
\]

\par
For a random variable $X$ and some constant $a \in \mathbb{R}$, the expectation of an indicator function is related to the CDF of $X$ as follows:
\[
\mathbb{E}[\mathbf{1}_{X \leq a}] = \int_{\Omega} \mathbf{1}_{\{X \leq a\}}\, dP = P(X \leq a) = F_X(a),
\]
where the indicator function of a set $A \in \mathcal{F}$ is defined by
\[
\mathbf{1}_A(\omega) = 
\begin{cases} 
1 & \omega \in A, \\
0 & \omega \notin A. 
\end{cases} \tag{1.1}
\]

\begin{definition}[Survival probability]
The survival probability is directly linked to the CDF. If $X$ is a random variable which denotes the lifetime, for example, in a population, then $\mathbb{P}[X \leq x]$ indicates the probability of not reaching the age $x$. The survival probability, defined by
\[
\mathbb{P}[X > x] = 1 - \mathbb{P}[X \leq x] = 1 - F_X(x),
\]
indicates the probability of surviving a lifetime of length $x$.
\end{definition}
\subsection{Lebesgue integral and convergence theorems}

The Lebesgue integral extends the classical Riemann integral by integrating with respect to a measure rather than over intervals. This allows integration of a wider class of functions and provides better convergence properties.

\par
One key difference is that the Lebesgue integral is defined by approximating functions via simple functions and measuring the size of their level sets, rather than partitioning the domain.

\par
The following convergence results are fundamental:

\begin{itemize}
    \item \textbf{Monotone Convergence Theorem (MCT):}  
    If $(X_n)$ is an increasing sequence of non-negative measurable functions, then
    \[
    \lim_{n \to \infty} \int X_n\, dP = \int \lim_{n \to \infty} X_n\, dP.
    \]

    \item \textbf{Fatou's Lemma:}  
    For a sequence $(X_n)$ of non-negative measurable functions,
    \[
    \int \liminf_{n \to \infty} X_n\, dP \leq \liminf_{n \to \infty} \int X_n\, dP.
    \]

    \item \textbf{Dominated Convergence Theorem (DCT):}  
    If $X_n \to X$ pointwise and $|X_n| \leq Y$ for some integrable function $Y$, then
    \[
    \lim_{n \to \infty} \int X_n\, dP = \int X\, dP.
    \]
\end{itemize}

\par
These results justify the interchange of limits and integrals under suitable conditions and are essential in probability theory and analysis.

\section{Filtration and Stochastic Processes}
\subsection{Stochastic processes and filtration}
We will often work with stochastic processes for the financial asset prices, and give some basic definitions for them here.
A \textbf{stochastic process}, $X(t)$, is a collection of random variables indexed by a time variable $t$.
Suppose we have a set of calendar dates/days, $T_1, T_2, \ldots, T_m$. Up to today, we have observed certain state values of the stochastic process $X(t)$. The past is known, and we therefore "see" the historical asset path. For the future we do not know the precise path but we may simulate the future according to some asset price distribution.

The mathematical tool which helps us describe the knowledge of a stochastic process up-to a certain time $T_i$ is the \textbf{sigma-field}, also known as \textbf{sigma-algebra}. The ordered sequence of sigma-fields is called a \textbf{filtration}, $\mathcal{F}(T_i) := \sigma(X(T_j) : 1 \leq j \leq i)$, generated by the sequence $X(T_j)$ for $1 \leq j \leq i$. The information available at time $T_i$ is thus described by a filtration. As we consider a sequence of observation dates, $T_1, \ldots, T_i$, we deal in fact with a sequence of filtrations, $\mathcal{F}(T_1) \subseteq \cdots \subseteq \mathcal{F}(T_i)$.

If we write that a process is $\mathcal{F}(T)$-measurable, we mean that at any time $t \leq T$, the realizations of this process are known. A simple example for this may be the market price of a stock and its historical values, i.e., we know the stock values up to today exactly, but we do not know any future values. We then say "the stock is today measurable". However, when we deal with an SDE model for the stock price, the value may be $T$ measurable, as we know the distribution for the period $T$ of a financial contract.

A stochastic process $X(t)$, $t \geq 0$, is said to be \textbf{adapted} to the filtration $\mathcal{F}(t)$, if
\[
\sigma(X(t)) \subseteq \mathcal{F}(t).
\]
By the term "adapted process" we mean that a stochastic process "cannot look into the future". In other words, for a stochastic process $X(t)$ its realizations (paths), $X(s)$ for $0 \leq s < t$, are known at time $s$ but not yet at time $t$.
\subsection{Martingales}

An important notion when dealing with stochastic processes is the martingale property.

\begin{definition}[Martingale]
Consider a probability space $(\Omega, \mathcal{F}, \mathbb{Q})$ equipped with a filtration $\{\mathcal{F}(t)\}_{t \ge 0}$.  
A right–continuous process $X(t)$ with left limits (a \emph{càdlàg} process) for $t \in [0,T]$ is said to be a martingale with respect to the filtration $\{\mathcal{F}(t)\}$ under the measure $\mathbb{Q}$ if, for all $s < t$, the following conditions hold:
\begin{align}
\mathbb{E}[\,|X(t)|\,] &< \infty, \\
\mathbb{E}[\,X(t)\mid \mathcal{F}(s)\,] &= X(s).
\end{align}
\end{definition}

The definition implies that the best prediction of a martingale’s future value is its present value. Indeed, for any $\Delta t > 0$,
\[
\mathbb{E}\big[X(t+\Delta t) - X(t)\mid \mathcal{F}(t)\big] = 0.
\]

\begin{remark}
A martingale represents a ``fair game'': once the current value $X(t)$ is known, the expected future change is zero. The process exhibits no predictable drift, and the conditional expectation of all future values equals the present value.
\end{remark}
\subsection{Adapted stochastic process}

We first need to introduce the binomial asset-pricing model in order to understand the following concepts.

\begin{definition}[Binomial asset pricing] 
In the \emph{binomial asset-pricing model}, time is divided into $N$ discrete steps.  
The price of a risky asset evolves according to
\[
S_{n+1} = 
\begin{cases}
u S_n, & \text{with probability } p,\\[4pt]
d S_n, & \text{with probability } 1-p,
\end{cases}
\]
where $u>1$ and $0<d<1$ are the up and down factors, respectively, and $S_0$ is the initial price.  
A risk-free asset grows deterministically at rate $r$, so its value satisfies
\[
B_{n+1} = (1+r) B_n.
\]
The model assumes $d < 1+r < u$ to exclude arbitrage.
\end{definition}

\begin{definition}
    Consider the binomial asset-pricing model and let
\[
M_0, M_1, \ldots, M_N
\]
be a sequence of random variables such that each $M_n$ depends only on the first $n$ coin tosses (with $M_0$ constant).  
Such a sequence is called an \emph{adapted stochastic process}.
\begin{enumerate}
    \item If for every $n = 0, 1, \ldots, N-1$,
    \[
    M_n = \mathbb{E}_n[M_{n+1}],
    \]
    we say that the process is a \emph{martingale}.

    \item If
    \[
    M_n \le \mathbb{E}_n[M_{n+1}], \qquad n = 0,1,\ldots,N-1,
    \]
    we say that the process is a \emph{submartingale}.

    \item If
    \[
    M_n \ge \mathbb{E}_n[M_{n+1}], \qquad n = 0,1,\ldots,N-1,
    \]
    we say that the process is a \emph{supermartingale}.
\end{enumerate}
\end{definition}

\textbf{Remark}  
The martingale property above is stated as a one-step-ahead condition,  
but it automatically implies an analogous condition for any number of steps.  
If $M_0, M_1, \ldots, M_N$ is a martingale and $n \le N-2$, then
\[
M_{n+1} = \mathbb{E}_{n+1}[M_{n+2}].
\]
Taking conditional expectations with respect to the information at time $n$, and using the tower property of conditional expectation, we obtain
\[
\mathbb{E}_n[M_{n+1}]
= \mathbb{E}_n\!\left( \mathbb{E}_{n+1}[M_{n+2}] \right)
= \mathbb{E}_n[M_{n+2}].
\]
Because $\mathbb{E}_n[M_{n+1}] = M_n$, this yields the “two-step-ahead’’ relation
\[
M_n = \mathbb{E}_n[M_{n+2}].
\]
Iterating the argument, for any $0 \le n \le m \le N$ we obtain the general multistep property:
\[
M_n = \mathbb{E}_n[M_m].
\]

\medskip

\textbf{Remark}  
A martingale has constant expectation over time.  
If $M_0, M_1, \ldots, M_N$ is a martingale, then
\[
M_0 = \mathbb{E}[M_n], \qquad n = 0,1,\ldots,N.
\]
Indeed, taking expectations on both sides of the martingale relation
\[
M_n = \mathbb{E}_n[M_{n+1}],
\]
and using the tower property, we find that $\mathbb{E}[M_n] = \mathbb{E}[M_{n+1}]$ for all $n$.  
Since $M_0$ is deterministic, $M_0 = \mathbb{E}[M_0]$, and the conclusion follows.
\subsection{Markov process}


\begin{definition}[Markov process] Consider the binomial asset-pricing model.  
Let $X_0, X_1, \ldots, X_N$ be an adapted process.  
If, for every $n$ between $0$ and $N-1$ and for every function $f(x)$,  
there exists another function $g(x)$ (depending on $n$ and $f$) such that
\begin{equation}
    \mathbb{E}_n[f(X_{n+1})] = g(X_n),
\end{equation}
we say that $X_0, X_1, \ldots, X_N$ is a Markov process.
\end{definition}

\begin{remark}
The Markov property expresses the idea that the process has no memory beyond its current state. Once $X_n$ is known, the conditional distribution of $X_{n+1}$ does not depend on the earlier values $X_0,\ldots,X_{n-1}$. Thus, the present state contains all the information necessary for predicting the next step.
\end{remark}


By definition, $\mathbb{E}_n[f(X_{n+1})]$ is a random quantity, since it depends on the first $n$ coin tosses. The Markov property states that this dependence occurs entirely through $X_n$; that is, all the information from the coin-toss history that is needed to evaluate $\mathbb{E}_n[f(X_{n+1})]$ is summarized by $X_n$.

At this stage, we are not focused on deriving an explicit formula for the function $g$, but rather on emphasizing its existence. The mere existence of such a function implies that if the payoff of a derivative security is random only through its dependence on $X_N$, then there is a version of the pricing algorithm that does not require storing the entire path of the process.

The martingale property arises as a particular instance of ${E}_n[f(X_{n+1})] = g(X_n)$ when  $f(x) = x$ and $g(x) = x$. For a process to qualify as Markov, it must satisfy ${E}_n[f(X_{n+1})] = g(X_n)$ for \emph{every} function $f$, with some corresponding function $g$.
Thus, while every martingale satisfies this relation for the specific choice $f(x) = x$, this alone does not ensure that the process is Markov. In fact, a general martingale need not be Markov.

Conversely, even when we restrict attention to the function $f(x) = x$, the Markov property merely requires that 
\[
\mathbb{E}_n[M_{n+1}] = g(M_n)
\]
for some function $g$; it does not insist that $g(x) = x$. Therefore, a Markov process is not necessarily a martingale.

\begin{definition}[Transition probabilities]
A Markov process $\{X_n\}$ is characterized by the transition probabilities
\[
P_n(x,A) := \mathbb{P}(X_{n+1} \in A \mid X_n = x),
\]
where $A$ is any measurable subset of the state space. For any function $f$, the function $g$ appearing in
\[
\mathbb{E}_n[f(X_{n+1})] = g(X_n)
\]
is given by the integral representation
\[
g(x) = \int f(y)\, P_n(x)dy.
\]
\end{definition}


The next lemma frequently provides a crucial step in confirming that a
given process possesses the Markov property.

\textbf{Lemma (Independence).}
In the $N$-period binomial asset–pricing model, let $n$ be an integer with $0 \leq n \leq N$. Suppose the random variables
\[
X_1,\ldots,X_K
\]
depend only on coin tosses $1$ through $n$, and the random variables
\[
Y_1,\ldots,Y_L
\]
depend only on coin tosses $n+1$ through $N$.  
(The superscripts $1,\ldots,K$ on $X$ and $1,\ldots,L$ on $Y$ are indices, not exponents.)

Let 
\[
f(x_1,\ldots,x_K,y_1,\ldots,y_L)
\]
be a function of dummy variables, and define the random variable
\begin{equation}
Z := f(X_1,\ldots,X_K,\,Y_1,\ldots,Y_L).
\end{equation}

Because the $X$'s depend only on the first $n$ coin tosses and the $Y$'s only on the remaining ones, the $X$'s are independent of the $Y$'s conditional on the information available at time $n$. Hence the conditional expectation satisfies
\begin{equation}
\mathbb{E}_n[Z]
  = \mathbb{E}_n\!\left[ f(X_1,\ldots,X_K,Y_1,\ldots,Y_L) \right]
  = g(X_1,\ldots,X_K),
\end{equation}
where the function $g$ is given by
\[
g(x_1,\ldots,x_K)
   := \mathbb{E}\!\left[ f(x_1,\ldots,x_K,Y_1,\ldots,Y_L) \right].
\]

Thus the conditional expectation with respect to $\mathcal{F}_n$ averages only over the randomness in $Y_1,\ldots,Y_L$, treating the $X$'s as fixed.

\section{Brownian Motion}
\subsection{Introduction}
The concept of \textit{Brownian motion} takes its origins in physics to describe the random movement of a particle suspended in a fluid. The mathematical formalization of this phenomenon, introduced by Wiener in the early 20th century, gave rise to one of the cornerstones of probability theory: the \textit{Wiener process} or \textit{standard Brownian motion}.

Brownian motion finds its particular significance in finance, where it serves as the paradigm of randomness in price movements.
\subsection{Definition and Fundamental Properties}

We first begin with the definition of standard Brownian motion. A stochastic process \( \{W_t\}_{t \geq 0} \) defined on a probability space \( (\Omega, \mathcal{F}, \mathbb{P}) \) is called a \textbf{standard Brownian motion} if it satisfies the following properties:
\begin{enumerate}
    \item \( W_0 = 0 \) almost surely;
    \[
    \mathcal{P}[W_U = 0] = 1.
    \]
    \item For every \( 0 \leq s < t \), the increment \( W_t - W_s \) is normally distributed with mean zero and variance \( t - s \), that is,
    \[
    W_t - W_s \sim \mathcal{N}(0, t-s);
    \]
    \item The increments are independent, meaning that for any sequence of times \( 0 \leq t_0 < t_1 < \cdots < t_n \),
    \[
    W_{t_1} - W_{t_0}, \, W_{t_2} - W_{t_1}, \, \ldots, \, W_{t_n} - W_{t_{n-1}}
    \quad \text{are independent};
    \]
    \item The sample paths of \( W_t \) are almost surely continuous.
\end{enumerate}


These properties uniquely characterize standard Brownian motion. The combination of Gaussian independent increments and continuity makes it a central process in modeling continuous random phenomena.
\subsection{Construction of Brownian Motion}
Up to this point, we have described the main properties that a Brownian motion should satisfy. 
However, stating its properties is not enough to guarantee that a process with such characteristics actually exists.
Hence, we need to prove the existence of such a process.
From a theoretical standpoint, the Brownian motion can be constructed in several ways, two classical and fundamental approaches are:

\begin{enumerate}
    \item \textbf{Kolmogorov's Construction}

    We begin by fixing a finite set of times $0 \le t_1 < t_2 < \dots < t_n$ and defining the joint distribution of the random vector
\[
(W_{t_1}, W_{t_2}, \dots, W_{t_n}).
\]
We require this vector to be multivariate normal with mean zero and covariance matrix
\[
\text{Cov}(W_{t_i}, W_{t_j}) = \min(t_i, t_j), \quad 1 \le i,j \le n.
\]
The intuition behind this covariance choice is as follows. For $s < t$, we can write
\[
W_t = W_s + (W_t - W_s).
\]
By definition of Brownian motion, the increment $(W_t - W_s)$ is independent of the past (including $W_s$) and has variance $t-s$. 
Therefore, the covariance satisfies
\[
\text{Cov}(W_s, W_t) = \text{Cov}(W_s, W_s + (W_t - W_s)) = \text{Var}(W_s) = s = \min(s,t),
\]
because the increment does not contribute to the covariance. 
In other words, the covariance function $\min(s,t)$ encodes two key properties at once: the independence of increments and the linear growth of variance with time.
    
    
Next, we must ensure that these finite-dimensional distributions are compatible with each other, a property called \emph{consistency}. 

A family of finite-dimensional distributions $\{\mu_{t_1, \dots, t_n}\}$ is \emph{consistent} if, for any subset of times, the marginal distribution of $\mu_{t_1, \dots, t_n}$ corresponding to the subset coincides with the distribution defined directly for those times. 
Intuitively, consistency ensures that all “pieces” of the process fit together.

Kolmogorov's Extension Theorem then guarantees the existence of a stochastic process with the given finite-dimensional distributions.

\begin{theorem}[Kolmogorov's Extension Theorem]
Let $\{\mu_{t_1, \dots, t_n}\}$ be a consistent family of finite-dimensional distributions indexed by finite ordered subsets of $[0, \infty)$. 
Then, there exists a probability space $(\Omega, \mathcal{F}, \mathbb{P})$ and a stochastic process $\{X_t\}_{t \ge 0}$ such that, for any $t_1 < \dots < t_n$, the joint distribution of $(X_{t_1}, \dots, X_{t_n})$ coincides with $\mu_{t_1, \dots, t_n}$.
\end{theorem}    


In our case, the Gaussian distributions defined above are consistent because taking the marginal of a multivariate normal distribution yields another multivariate normal with the correct mean and covariance. 
Hence, Kolmogorov's theorem ensures the existence of a stochastic process $(W_t)_{t \ge 0}$ with the desired finite-dimensional distributions.

While existence is guaranteed, continuity of sample paths is not. To obtain continuity, we use the \textit{Kolmogorov–Chentsov Continuity Theorem}.

\begin{theorem}[Kolmogorov–Chentsov Continuity Theorem]
Let $\{X_t\}_{t \ge 0}$ be a stochastic process. 
If there exist constants $\alpha, \beta, C > 0$ such that
\[
\mathbb{E}[|X_t - X_s|^{\alpha}] \le C |t-s|^{1+\beta}, \quad \forall s,t \ge 0,
\]
then there exists a \emph{modification} of $X_t$ with continuous sample paths almost surely.
\end{theorem}

For Brownian motion, the fourth moment of the increments satisfies
\[
\mathbb{E}[(W_t - W_s)^4] = 3 (t - s)^2,
\]
which fits the theorem with $\alpha = 4$ and $\beta = 1$. 
Therefore, a continuous version of the process exists.

In conclusion, we have constructed a stochastic process $(W_t)_{t \ge 0}$ that satisfies:

\begin{itemize}
    \item $W_0 = 0$ almost surely;
    \item independent Gaussian increments with variance equal to the length of the interval;
    \item covariance function $\text{Cov}(W_s, W_t) = \min(s,t)$, encoding both independence of increments and linear variance growth;
    \item continuous sample paths almost surely.
\end{itemize}

    
    \item \textbf{The Random Walk Limit}
    
Another method used to create the classic Brownian Motion is the Random Walk Limit. This method provides an intuitive link between discrete-time models and the continuous-time stochastic dynamics used in finance. It considers a sequence of independent and identically distributed (i.i.d.) random variables $\{\xi_i\}_{i \in \mathbb{N}}$, with zero mean ($E[\xi_i] = 0$) and unit variance ($\text{Var}(\xi_i) = 1$). The rescaled random walk process $\{W^{(n)}_t\}_{t \geq 0}$ is defined as:
    
    $$
    W^{(n)}_t = \frac{1}{\sqrt{n}} \sum_{i=1}^{\lfloor nt \rfloor} \xi_i
    $$
    
where $\lfloor nt \rfloor$ denotes the integer part of $nt$, corresponding to the number of steps up to time $t$.

The scaling factor $\frac{1}{\sqrt{n}}$ is crucial: it ensures that, as the number of steps $n$ increases, the variance of $W^{(n)}_t$ grows linearly with time, matching the defining property of Brownian motion. 

As $n \to \infty$, the process $W^{(n)}_t$ converges in distribution to a continuous-time process $\{W_t\}_{t \ge 0}$, which is the standard Brownian motion. 

Intuitively, as the time steps become finer and the step sizes are appropriately scaled, the discrete jumps of the random walk become infinitesimally small and occur infinitely often in a finite interval. In the limit, the random walk “smooths out” into a continuous, nowhere-differentiable path. The key features of Brownian motion: zero mean, independent increments, variance proportional to elapsed time and continuity of paths emerge naturally from this limiting procedure. 
\end{enumerate}

% Asset/code block omitted pending provenance acceptance.
\subsection{Basic Properties of Brownian Motion}

Based on the definition provided in the previous section, the standard Brownian motion $\{W_t\}_{t \ge 0}$ possesses several fundamental properties that are crucial for its application in stochastic calculus and financial modeling.
\paragraph{Expected Value and Covariance}
The Gaussian nature and the property of zero initial state ($W_0=0$) lead directly to the mean and covariance functions:

\begin{enumerate}
    \item \textbf{Expected Value (Mean):} The expected value of Brownian motion at any time $t$ is zero.
    \[
    \mathbb{E}[W_t] = 0, \quad \forall t \ge 0.
    \]
    This follows from the fact that $W_t = (W_t - W_0)$ is normally distributed as $\mathcal{N}(0, t)$, which has a mean of zero.

    \item \textbf{Variance:} The variance of Brownian motion is linear in time.
    \[
    \text{Var}(W_t) = \mathbb{E}[W_t^2] = t, \quad \forall t \ge 0.
    \]

    \item \textbf{Covariance Function:} For any $s, t \ge 0$, the covariance between $W_s$ and $W_t$ is given by the minimum of the two times:
    \[
    \text{Cov}(W_s, W_t) = \min(s, t).
    \]
    If we assume, without loss of generality, that $s \le t$, we can derive this property as follows, leveraging the independence of increments:
    \begin{align*}
    \text{Cov}(W_s, W_t) &= \text{Cov}(W_s, W_s + (W_t - W_s)) \\
    &= \text{Cov}(W_s, W_s) + \text{Cov}(W_s, W_t - W_s) \\
    &= \text{Var}(W_s) + 0 \\
    &= s = \min(s, t).
    \end{align*}
\end{enumerate}

% Asset/code block omitted pending provenance acceptance.
\paragraph{Martingale Property}
The independence of future increments from the past information is the key feature that establishes the Brownian motion as a martingale.


For standard Brownian motion $\{W_t\}_{t \ge 0}$ and its natural filtration $\{\mathcal{F}_t\}$, we have:
\begin{align*}
\mathbb{E}[W_t | \mathcal{F}_s] &= \mathbb{E}[W_s + (W_t - W_s) | \mathcal{F}_s] \\
&= \mathbb{E}[W_s | \mathcal{F}_s] + \mathbb{E}[W_t - W_s | \mathcal{F}_s] \\
&= W_s + \mathbb{E}[W_t - W_s] \\
&= W_s + 0 = W_s.
\end{align*}
The steps rely on the facts that $W_s$ is $\mathcal{F}_s$-measurable, and the increment $(W_t - W_s)$ is independent of $\mathcal{F}_s$ and has zero mean. Thus, $\mathbf{W_t}$ is a martingale.
\paragraph{Markov Property}
Brownian motion is also a Markov Process. 
For Brownian motion, $W_t$ can be written as $W_t = W_s + (W_t - W_s)$. Since the future increment $(W_t - W_s)$ is independent of the past $\{\mathcal{F}_s\}$ and $W_s$ summarizes all the necessary history for the transition, the future movement starting at $W_s$ is independent of how $W_s$ was reached. This makes $\mathbf{W_t}$ a Markov process.
\paragraph{Exponential Transformation and Martingale (Girsanov's Theorem)}
A particularly important transformation in mathematical finance is the exponential of the Brownian motion, especially the process known as the exponential martingale or Doléans-Dade exponential, which appears in the context of the Girsanov Theorem.

The process $Z_t$ defined by:
\[
Z_t = \exp \left( \alpha W_t - \frac{1}{2} \alpha^2 t \right)
\]
is a martingale for any real constant $\alpha$. This process, often called the Novikov Martingale (or simply the exponential martingale), is a key component in changing the probability measure (change of measure). The term $-\frac{1}{2}\alpha^2 t$ acts as a compensator to ensure that $\mathbb{E}[Z_t] = 1$, which is a necessary condition for a non-negative martingale with $Z_0=1$.

We confirm the martingale property using the expectation $\mathbb{E}[Z_t | \mathcal{F}_s] = Z_s$:
\[
\mathbb{E}[Z_t | \mathcal{F}_s] = Z_s \quad \text{almost surely}.
\]
This property will be fundamental later in deriving the risk-neutral measure in the Black-Scholes model.
\subsection{Filtration for Brownian Motion}

When applied to Brownian motion $\{W_t\}_{t \ge 0}$, the filtration $\{\mathcal{F}_t\}_{t \ge 0}$ represents the increasing collection of information generated by the process up to time $t$. This filtration must satisfy the following key properties reflecting the nature of Brownian motion:

\begin{enumerate}
    \item \textbf{Information accumulation:} For all $0 \le s < t$, the sigma-algebra $\mathcal{F}_s$ is contained in $\mathcal{F}_t$, meaning that information available at an earlier time is also available at later times.
    
    \item \textbf{Adaptedness:} The Brownian motion is adapted to $\{\mathcal{F}_t\}$, i.e., for each $t \ge 0$, the random variable $W_t$ is measurable with respect to $\mathcal{F}_t$. Thus, the value of $W_t$ is completely determined by the information known up to time $t$.
    
    \item \textbf{Independence of future increments:} For any $0 \le t < u$, the increment $W_u - W_t$ is independent of the sigma-algebra $\mathcal{F}_t$. This expresses the fundamental property that the future evolution of the Brownian motion cannot be predicted from past information.
\end{enumerate}

These properties formalize the intuitive understanding that Brownian motion “reveals” information progressively and that its future increments behave like fresh randomness independent of the past. This independence underlies many important theoretical results and practical models, including the efficient market hypothesis in financial mathematics, where past information cannot be used to forecast future price changes modeled by Brownian motion.

Moreover, any stochastic process $\{\Delta_t\}_{t \ge 0}$ is said to be adapted to the filtration $\{\mathcal{F}_t\}$ if, for every $t$, $\Delta_t$ is measurable with respect to $\mathcal{F}_t$. This ensures that $\Delta_t$ depends only on information available up to time $t$, a crucial requirement in stochastic calculus and modeling.
\section{SDEs and Geometric Brownian Motion}
In the world of finacial markets, asset's prices reflect both the present value of the firm and expectations about their future performance. These expectations, however, fluctuate over time introducing an uncertainty component into the asset price. The resulting randomness in asset prices is possible to model through the use of stochastic differential equations (SDEs).
SDEs provide analytical tools for validating numerical schemes and serve as building blocks for more sophisticated stochastic models of market dynamics.
Among all continuous-time models, the most fundamental is the Geometric Brownian Motion (GBM). In this framework, the logarithm of the asset price follows an arithmetic Brownian motion driven by a Wiener process. This ensures properties that align well with real-world financial data: prices remain strictly positive, returns are continuous and volatility acts proportionally to the current price level.
However, formulating and analysing the GBM model rigorously requires the mathematical machinery of stochastic calculus. Before introducing GBM explicitly, we therefore present the notions of stochastic integrands and square-integrable processes, develop the Itô integral, and establish the general theory of SDEs.
\subsection{Stochastic Integrands and SDEs}
The Stochastic Differential Equation (SDE) is defined in terms of the Itô integral with respect to the standard Brownian motion $W_t$. To ensure the existence and properties of this integral, the functions we integrate (the integrands) must satisfy specific mathematical conditions.
\paragraph{Stochastic Integrands and Square-Integrable Variables}
The Itô integral, $\int_0^T \Phi_t dW_t$, is constructed as a limit of sums (analogous to the Riemann sum) built on the properties of martingales and $L^2$ theory. The class of processes that can be integrated with respect to $W_t$ is restricted to ensure convergence in the mean-square sense.

\begin{definition}[Admissible Integrand]
A stochastic process $\Phi = \{\Phi_t\}_{t \geq 0}$ defined on the probability space $(\Omega, \mathcal{F}, \mathbb{P})$ is classified as an admissible integrand (or belonging to the space $\mathcal{L}^2(0, T)$) if it satisfies the following two conditions:
\begin{enumerate}
    \item Adaptedness: $\Phi_t$ is $\mathcal{F}_t$-adapted, meaning $\Phi_t$ depends only on the information available up to time $t$.
    \item Square-Integrability: $\Phi_t$ is square-integrable over the time interval $[0, T]$, i.e.,
    \[
    \mathbb{E} \left[ \int_0^T \Phi_t^2 dt \right] < \infty, \quad \text{for all } T > 0.
    \]
\end{enumerate}
\end{definition}
The condition of square-integrability for the process ensures that the variance of the resulting Itô integral, $\text{Var}\left( \int_0^T \Phi_t dW_t \right) = \mathbb{E} \left[ \int_0^T \Phi_t^2 dt \right]$, remains finite. Similarly, a random variable $X$ is square-integrable if its second moment is finite: $\mathbb{E}[X^2] < \infty$.
\paragraph{Stochastic Differential Equations (SDEs)}
A Stochastic Differential Equation describes a continuous-time process whose evolution is partially deterministic and partially stochastic, driven by Brownian motion.

\begin{definition}[General SDE]
A general form of a one-dimensional Stochastic Differential Equation is an equation for the stochastic process $\{X_t\}_{t \ge 0}$ written as:
\[
dX_t = \mu(t, X_t) dt + \sigma(t, X_t) dW_t, \quad X_0 = x_0,
\]
which is interpreted in its integral form:
\[
X_t = x_0 + \int_0^t \mu(u, X_u) du + \int_0^t \sigma(u, X_u) dW_u.
\]
\end{definition}
Here:
\begin{itemize}
    \item $X_t$ is the unknown stochastic process, often modeling an asset price or interest rate.
    \item $\mu(t, X_t)$ is the drift coefficient (or function), representing the deterministic component of the change, related to the expected growth.
    \item $\sigma(t, X_t)$ is the diffusion coefficient (or function), representing the intensity of the noise component, proportional to the volatility.
    \item The first integral is a standard Lebesgue integral, while the second is the Itô stochastic integral, which requires the integrand $\sigma(u, X_u)$ to be admissible.
\end{itemize}
Existence and uniqueness of the solution $X_t$ are typically guaranteed if the coefficients $\mu$ and $\sigma$ satisfy the local Lipschitz and linear growth conditions.
\subsection{Itô’s Formula}

In standard differential calculus, the chain rule is used to find the differential of a composite function, $Y_t = f(t, X_t)$, where $X_t$ is a deterministic differentiable function. However, the standard chain rule does not apply when $X_t$ is an Itô process.

The failure of the classical chain rule in the stochastic setting is fundamentally due to the \textbf{non-zero quadratic variation} of Brownian motion. For a smooth deterministic function, $(dx)^2 \approx 0$ as $dt \to 0$. However, for Brownian motion $W_t$, the differential $(dW_t)^2$ is of order $dt$ (formally, $E[(dW_t)^2] = dt$). When performing a Taylor expansion of $f(t, X_t)$, the second-order term $(dX_t)^2$ does not vanish in the limit, but instead contributes a first-order term to the differential. This "extra term" is known as the Itô term or the drift correction.
\paragraph{Construction via Summation}
The Itô formula can be heuristically derived by considering a partition of the time interval $[0, T]$, where $0 = t_0 < t_1 < \dots < t_n = T$. Let $\Delta X_j = X_{t_{j+1}} - X_{t_j}$. Using a Taylor expansion:
\[
f(T, X_T) - f(0, X_0) = \sum_{j=0}^{n-1} [f(t_{j+1}, X_{t_{j+1}}) - f(t_j, X_{t_j})]
\]
Expanding the terms inside the sum:
\[
\Delta f \approx \sum \frac{\partial f}{\partial t} \Delta t_j + \sum \frac{\partial f}{\partial x} \Delta X_j + \frac{1}{2} \sum \frac{\partial^2 f}{\partial x^2} (\Delta X_j)^2 + \dots
\]
As the mesh size of the partition goes to zero, the term $\sum (\Delta X_j)^2$ converges in probability to the quadratic variation $\langle X \rangle_T$. This summation process highlights that the Itô integral is defined by evaluating the integrand at the \textbf{left endpoint} of each sub-interval, i.e., $f(t_j, X_{t_j})$.

\begin{theorem}[Itô’s Formula]
Let $X_t$ be an Itô process defined by: $dX_t = \mu_t dt + \sigma_t dW_t$.
For $f \in C^{1,2}([0, \infty) \times \mathbb{R})$, the differential $dY_t$ of $Y_t = f(t, X_t)$ is:
\[
dY_t = \left( \frac{\partial f}{\partial t} + \frac{\partial f}{\partial x} \mu_t + \frac{1}{2} \frac{\partial^2 f}{\partial x^2} \sigma_t^2 \right) dt + \frac{\partial f}{\partial x} \sigma_t dW_t.
\]
\end{theorem}
\paragraph{Comparison: Itô vs. Stratonovich}
The choice of the evaluation point in the Riemann-sum construction leads to different stochastic integrals:

\begin{itemize}
    \item \textbf{Itô Integral ($\int X_t dW_t$):} Evaluates the integrand at the \textbf{left endpoint} ($t_j$). It is a martingale (under certain conditions) and is "non-anticipating," making it ideal for finance as it respects causality (information available at time $t$). However, it does not follow classical calculus rules.
    \item \textbf{Stratonovich Integral ($\int X_t \circ dW_t$):} Evaluates the integrand at the \textbf{midpoint} ($ \frac{t_j + t_{j+1}}{2} $).
     It \textbf{satisfies the classical chain rule} (the second-order term disappears in its formulation). However, it is not generally a martingale and "looks into the future" within each small interval, making it less intuitive for modeling financial price processes.
    
\end{itemize}
\subsection{Geometric Brownian Motion}

We now apply the general theory of stochastic integrals and Itô calculus to derive and analyse the most fundamental model for asset prices in continuous time: the \emph{Geometric Brownian Motion} (GBM). This model captures proportional growth, multiplicative noise and strict positivity of prices, and forms the basis of the Black--Scholes option pricing framework.
\paragraph{Definition of the GBM SDE}

\begin{definition}[Geometric Brownian Motion]
A stochastic process $S = \{S_t\}_{t \ge 0}$ is said to follow a \emph{Geometric Brownian Motion} if it satisfies the Stochastic Differential Equation
\begin{equation}
dS_t = \mu S_t\, dt + \sigma S_t\, dW_t, \qquad S_0 > 0,
\end{equation}
where
\begin{itemize}
    \item $\mu \in \mathbb{R}$ is the \emph{drift}, representing the expected rate of return,
    \item $\sigma > 0$ is the \emph{volatility}, representing the intensity of stochastic fluctuations,
    \item $W_t$ is a standard Brownian motion on the filtered probability space.
\end{itemize}
\end{definition}

The coefficients of the SDE,
\[
b(t,s) = \mu s, \qquad a(t,s) = \sigma s,
\]
are globally Lipschitz and satisfy linear growth conditions. Thus, by the fundamental existence and uniqueness theorem for SDEs, equation \sourceref{eq:GBM-SDE} has a unique strong solution.
\paragraph{Logarithmic Dynamics via Itô's Formula}

A key observation is that the logarithm of $S_t$ evolves as an arithmetic Brownian motion.  
Let $Y_t = \ln(S_t)$. Applying Itô’s formula to the function $f(s) = \ln s$ yields
\[
df(S_t)
= \frac{1}{S_t} dS_t 
- \frac12 \frac{1}{S_t^2} (dS_t)^2.
\]
Since $(dS_t)^2 = \sigma^2 S_t^2 dt$, substituting the SDE \sourceref{eq:GBM-SDE} gives
\begin{equation}
dY_t = \left( \mu - \tfrac{1}{2}\sigma^2 \right) dt + \sigma\, dW_t.
\end{equation}
Thus the logarithmic price process follows a linear SDE with constant coefficients.
\paragraph{Explicit Solution and Log-Normality}

Integrating \sourceref{eq:logS-dynamics} from $0$ to $t$ yields
\[
\ln(S_t)
= \ln(S_0)
+ \left(\mu - \tfrac{1}{2}\sigma^2\right)t
+ \sigma W_t.
\]
Exponentiating gives the explicit closed-form solution:
\begin{equation}
S_t
= S_0 \exp\!\left[
     \left(\mu - \tfrac12\sigma^2\right)t
     + \sigma W_t
   \right].
\end{equation}

From this expression we obtain several fundamental properties:

\begin{enumerate}
    \item \textbf{Strict positivity}:
    \[
    S_t > 0 \quad \text{a.s. for all } t \ge 0.
    \]

    \item \textbf{Log-normal distribution}:  
    Since $W_t \sim \mathcal{N}(0,t)$, it follows that
    \[
    \ln(S_t)
    \sim \mathcal{N}\!\left(
        \ln(S_0) + \left(\mu - \tfrac12 \sigma^2\right)t,
        \; \sigma^2 t
    \right).
    \]

    \item \textbf{Moments}:
    \[
    \mathbb{E}[S_t] = S_0 e^{\mu t}, \qquad
    \mathrm{Var}(S_t) = S_0^2 e^{2\mu t} \big( e^{\sigma^2 t} - 1 \big).
    \]
\end{enumerate}


% Asset/code block omitted pending provenance acceptance.
\paragraph{Interpretation in Financial Modelling}

The GBM SDE captures two essential facts commonly observed in financial markets, making it a natural starting point for modelling asset prices in continuous time.

\begin{itemize}
    \item \textbf{Proportional volatility:}  
    The diffusion term $\sigma S_t dW_t$ is proportional to the current level of the asset price. This implies that the magnitude of random fluctuations increases with the price itself. As a consequence, the model assumes that returns, rather than absolute price changes, exhibit a constant level of volatility over time. This property is consistent with empirical observations in financial data, where percentage changes tend to be more stable than absolute variations.

    \item \textbf{Strict positivity of prices:}  
    The explicit solution of the GBM, given in \sourceref{eq:GBM-solution}, is expressed as an exponential function of a Brownian motion. Since the exponential function is always strictly positive, it follows that $S_t > 0$ almost surely for all $t \ge 0$. This property is particularly important in financial modelling, as it ensures that asset prices cannot become negative.
\end{itemize}

\section{Risk-Neutral Pricing}
\subsection{Asset Prices Under the Physical Measure}

Throughout this chapter we assume that the asset price process satisfies
\[
\mathbb{E}[\,|S(t)|\,] < \infty,
\]
so that expectations of payoffs are well defined.

Under the real-world probability measure $\mathbb{P}$, the dynamics of a stock price are typically modelled by a geometric Brownian motion
\[
dS(t) = \mu S(t)\,dt + \sigma S(t)\, dW^{\mathbb{P}}(t),
\]
where $\mu$ denotes the expected rate of return and $\sigma>0$ the volatility.

The explicit solution of this stochastic differential equation is
\[
S(t) = S(t_0)\exp\left[\left(\mu-\frac{1}{2}\sigma^2\right)(t-t_0) + \sigma\big(W^{\mathbb P}(t)-W^{\mathbb P}(t_0)\big)\right].
\]

Taking conditional expectations yields
\[
\mathbb{E}_{\mathbb{P}}\!\left[ S(t)\mid \mathcal{F}_{t_0} \right]
= S(t_0)e^{\mu(t-t_0)}.
\]

Hence, over a short interval $\Delta t$,
\[
\mathbb{E}_{\mathbb{P}}\!\left[ S(t+\Delta t)\mid \mathcal{F}_t \right]
= S(t)e^{\mu\Delta t}.
\]

Since typically $\mu \neq r$, the discounted price process
\[
\tilde S(t)=e^{-rt}S(t)
\]
satisfies
\[
d\tilde S(t)
=
(\mu-r)\tilde S(t)dt
+
\sigma\tilde S(t)dW^{\mathbb P}(t),
\]
and therefore it is generally \emph{not} a martingale under the physical measure $\mathbb{P}$.
\subsection{Risk-Neutral Measure and Discounted Martingales}

A central idea of modern asset pricing theory is to introduce an alternative probability measure $\mathbb{Q}$, called the \emph{risk-neutral measure}, under which discounted asset prices become martingales.

\par
The passage from $\mathbb{P}$ to $\mathbb{Q}$ is achieved via a change of measure. Define the process
\[
\lambda := \frac{\mu - r}{\sigma},
\]
called the \textit{market price of risk}. Then, by \textit{Girsanov's theorem}, there exists a probability measure $\mathbb{Q}$ equivalent to $\mathbb{P}$ such that the process
\[
W^{\mathbb{Q}}(t) := W^{\mathbb{P}}(t) + \lambda t
\]
is a Brownian motion under $\mathbb{Q}$.

\par
Equivalently, the change of measure can be expressed via the Radon--Nikodym derivative
\[
\frac{d\mathbb{Q}}{d\mathbb{P}}\Big|_{\mathcal{F}_t}
=
\exp\left(
- \lambda W^{\mathbb{P}}(t)
- \frac{1}{2}\lambda^2 t
\right).
\]

\par
Under $\mathbb{Q}$ the asset price dynamics become
\[
dS(t) = r S(t)\,dt + \sigma S(t)\, dW^{\mathbb{Q}}(t),
\]
where $r$ is the risk-free interest rate.

The money market account
\[
M(t) = e^{r(t-t_0)}, \qquad dM(t)=rM(t)\,dt,
\]
serves as the numéraire.

Defining the discounted stock price
\[
\tilde S(t)=\frac{S(t)}{M(t)}=e^{-rt}S(t),
\]
Itô's lemma yields
\[
d\tilde S(t)=\sigma\tilde S(t)dW^{\mathbb Q}(t),
\]
which implies that $\tilde S(t)$ is a martingale under $\mathbb Q$. In particular,
\[
\mathbb{E}_{\mathbb{Q}}
\!\left[
\frac{S(t+\Delta t)}{M(t+\Delta t)}
\middle|
\mathcal{F}_t
\right]
=
\frac{S(t)}{M(t)}.
\]

\par
This property characterizes the risk-neutral measure. More generally, according to the \emph{Fundamental Theorem of Asset Pricing}, the absence of arbitrage in a financial market is equivalent to the existence of an \emph{Equivalent Martingale Measure} (EMM) under which discounted asset prices are martingales.
\subsection{The Non-Uniqueness of the Equivalent Martingale Measure in the Lévy Framework}

In models where the stock price dynamics include jumps, such as those driven by Lévy processes, the financial market is typically \emph{incomplete}. In incomplete markets perfect replication of derivative payoffs is generally not possible.

As a consequence, the risk-neutral measure is no longer unique; instead, there exist infinitely many equivalent martingale measures.

A common approach in Lévy models is therefore to select a convenient martingale measure and compute derivative prices as discounted expectations under this measure.
\paragraph{The Equivalent Martingale Measure (EMM) Condition}

An Equivalent Martingale Measure (EMM), denoted by $\mathbb Q$, must satisfy two conditions:

\begin{enumerate}
\item $\mathbb Q$ is \emph{equivalent} to the physical measure $\mathbb P$, meaning that they share the same null sets.
\item The discounted stock price process $e^{-rt}S(t)$ is a martingale under $\mathbb Q$.
\end{enumerate}

The martingale property implies
\[
\mathbb E^{\mathbb Q}[e^{-rt}S(t)\mid\mathcal F_0]=S_0,
\]
which can be rewritten as
\[
\mathbb E^{\mathbb Q}[S(t)\mid\mathcal F_0]=S_0e^{rt}.
\]
\paragraph{Deriving the Martingale Constraint for the Lévy Exponent}

Consider a stock price process modeled as a geometric Lévy process
\[
S(t)=S_0e^{X_L(t)},
\]
where $X_L(t)$ is a Lévy process with $X_L(0)=0$.

Taking expectations gives
\[
\mathbb E^{\mathbb Q}[S(t)\mid\mathcal F_0]
=
S_0\,\mathbb E^{\mathbb Q}\!\left[e^{X_L(t)}\right].
\]

The martingale condition therefore requires
\[
\mathbb E^{\mathbb Q}\!\left[e^{X_L(t)}\right]=e^{rt}.
\]

The Lévy–Khintchine representation gives the characteristic function
\[
\phi_{X_L}(u)
=
\mathbb E^{\mathbb Q}\!\left[e^{iuX_L(t)}\right]
=
\exp\{t\psi_Q(u)\},
\]
where the characteristic exponent is
\[
\psi_Q(u)
=
i\mu_Lu
-
\frac{\sigma_L^2}{2}u^2
+
\int_{\mathbb R}
\left(e^{iux}-1-iux\mathbf 1_{|x|\le1}\right)F_L(dx).
\]

Evaluating the characteristic function at $u=-i$ yields
\[
\mathbb E^{\mathbb Q}\!\left[e^{X_L(t)}\right]
=
\exp\left\{
t\left(
\mu_L
+
\frac{\sigma_L^2}{2}
+
\int_{\mathbb R}
\left(e^x-1-x\mathbf 1_{|x|\le1}\right)F_L(dx)
\right)
\right\}.
\]

The martingale condition therefore implies the constraint
\[
\mu_L
+
\frac{\sigma_L^2}{2}
+
\int_{\mathbb R}
\left(e^x-1-x\mathbf 1_{|x|\le1}\right)F_L(dx)
=
r.
\]

This equation must be satisfied by any probability measure $\mathbb Q$ that serves as an equivalent martingale measure. Since multiple parameter combinations may satisfy this constraint, the EMM is generally not unique in Lévy models.
\subsection{Risk-Neutral Valuation in the Binomial Model}

In the $N$-period binomial asset pricing model, derivative pricing relies on the risk-neutral probability measure $\mathbb Q$.

The absence of arbitrage requires
\[
d < 1+r < u,
\]
where $u$ and $d$ are the up and down factors of the stock price.

Under this condition the risk-neutral probability is uniquely determined by
\[
q=\frac{(1+r)-d}{u-d}.
\]
\paragraph{The Martingale Property}

Under the risk-neutral measure $\mathbb Q$, discounted price processes are martingales.

If $V_n$ denotes the price of a derivative security at time $n$, the martingale property implies
\[
\frac{V_n}{(1+r)^n}
=
\mathbb E_{\mathbb Q}
\left[
\frac{V_{n+1}}{(1+r)^{n+1}}
\middle|
\mathcal F_n
\right],
\qquad n=0,\dots,N-1.
\]
\paragraph{The Risk-Neutral Pricing Formula}

For a derivative that pays $V_N$ at time $N$, the arbitrage-free price at time $n$ is given by
\[
V_n
=
\mathbb E_{\mathbb Q}
\left[
\frac{V_N}{(1+r)^{N-n}}
\middle|
\mathcal F_n
\right].
\]

Thus the price of the derivative is the conditional expectation of its discounted payoff under the risk-neutral measure.

\section{The Black--Scholes--Merton Model}

The Black--Scholes--Merton framework represents the analytical pinnacle of the stochastic calculus developed in the preceding sections. This chapter provides a formal derivation of the pricing formula for a European call option, utilizing the Risk-Neutral Valuation principle.
\subsection{Derivatives}
Before introducing the pricing framework, it is useful to recall the nature of the financial instruments under consideration. A \textit{derivative} is a financial contract whose value depends on the price of an underlying asset, such as a stock. Among the most common derivatives are \textit{options}, which grant the holder the right, but not the obligation, to buy or sell the underlying asset at a predetermined price $K$ (the \textit{strike}) at a future time $T$ (the \textit{maturity}).

A \textit{European call option} gives the holder the right to buy the underlying asset at price $K$ at maturity $T$. The payoff of the contract at time $T$ is therefore

\[
\max(S_T - K, 0)
\]

which reflects the fact that the holder exercises the option only when the market price $S_T$ exceeds the strike price.

Conversely, a \textit{European put option} gives the right to sell the asset at price $K$ at maturity $T$. Its payoff is

\[
\max(K - S_T, 0)
\]

which becomes positive only when the asset price falls below the strike price.

Option pricing models, such as the Black--Scholes--Merton framework, aim to determine the fair present value of these contingent payoffs under suitable probabilistic assumptions on the dynamics of the underlying asset.
\subsection{Theoretical Assumptions}
The model operates under a probability space $(\Omega, \mathcal{F}, \mathbb{Q})$ where the following conditions hold:
\begin{enumerate}
    \item The stock price $S_t$ follows a Geometric Brownian Motion (GBM): $dS_t = r S_t dt + \sigma S_t dW^{\mathbb{Q}}_t$.
    \item The risk-free rate $r$ and volatility $\sigma$ are constant.
    \item Markets are frictionless (no transaction costs, continuous trading).
    \item There are no arbitrage opportunities.
\end{enumerate}
\subsection{Full Derivation of the Pricing Formula}

Let $C(S, t)$ be the price of a European call option with strike $K$ and maturity $T$. According to the \textit{Risk-Neutral Valuation Thesis}, the price is the discounted expected value of the payoff under the equivalent martingale measure $\mathbb{Q}$:
\begin{equation}
C(S, t) = e^{-r\tau} \mathbb{E}^{\mathbb{Q}} \left[ \max(S_T - K, 0) \mid \mathcal{F}_t \right]
\end{equation}
where $\tau = T-t$. From the explicit solution of the GBM, we know:
\[ S_T = S_t \exp\left( \left(r - \frac{1}{2}\sigma^2\right)\tau + \sigma\sqrt{\tau}Z \right), \quad Z \sim \mathcal{N}(0,1) \]

To evaluate the expectation, we integrate with respect to the standard normal density $\phi(z) = \frac{1}{\sqrt{2\pi}}e^{-z^2/2}$:
\[ C(S, t) = e^{-r\tau} \int_{-\infty}^{\infty} \max\left( S_t e^{(r - \frac{1}{2}\sigma^2)\tau + \sigma\sqrt{\tau}z} - K, 0 \right) \phi(z) dz \]

The payoff is non-zero only when $S_T > K$. Solving for $z$:
\[ \ln(S_t) + (r - \frac{1}{2}\sigma^2)\tau + \sigma\sqrt{\tau}z > \ln(K) \implies z > \frac{\ln(K/S_t) - (r - \frac{1}{2}\sigma^2)\tau}{\sigma\sqrt{\tau}} \]
Let $d_2 = \frac{\ln(S_t/K) + (r - \frac{1}{2}\sigma^2)\tau}{\sigma\sqrt{\tau}}$. The condition becomes $z > -d_2$. We split the integral:
\[ C = e^{-r\tau} \int_{-d_2}^{\infty} S_t e^{(r - \frac{1}{2}\sigma^2)\tau + \sigma\sqrt{\tau}z} \phi(z) dz - e^{-r\tau} \int_{-d_2}^{\infty} K \phi(z) dz \]

\textbf{First Integral ($I_1$):}
\[ I_1 = S_t \int_{-d_2}^{\infty} \frac{1}{\sqrt{2\pi}} e^{-\frac{1}{2}\sigma^2\tau + \sigma\sqrt{\tau}z - \frac{1}{2}z^2} dz = S_t \int_{-d_2}^{\infty} \frac{1}{\sqrt{2\pi}} e^{-\frac{1}{2}(z - \sigma\sqrt{\tau})^2} dz \]
Substituting $y = z - \sigma\sqrt{\tau}$, and noting that as $z \to -d_2$, $y \to -(d_2 + \sigma\sqrt{\tau}) = -d_1$:
\[ I_1 = S_t \int_{-d_1}^{\infty} \phi(y) dy = S_t N(d_1) \]

\textbf{Second Integral ($I_2$):}
\[ I_2 = K e^{-r\tau} \int_{-d_2}^{\infty} \phi(z) dz = K e^{-r\tau} N(d_2) \]

Combining $I_1$ and $I_2$, we obtain the \textbf{Black--Scholes Formula}:

\[ C(S, t) = S_t N(d_1) - K e^{-r\tau} N(d_2) \]
where:
\[ d_1 = \frac{\ln(S_t/K) + (r + \frac{1}{2}\sigma^2)\tau}{\sigma\sqrt{\tau}}, \quad d_2 = d_1 - \sigma\sqrt{\tau} \]
\subsection{Small introduction to Greeks}

The Greeks represent the sensitivities of the option price to the underlying parameters. They are fundamental for the dynamic replication strategy (Delta hedging).

% Asset/code block omitted pending provenance acceptance.




\begin{remark}
The term $N(d_2)$ represents the \textbf{exercise probability} under the risk-neutral measure, while $N(d_1)$ is the \textbf{hedge ratio}. The discrepancy between these two terms is a direct consequence of the volatility-induced drift correction required in stochastic environments.
\end{remark}
